\section{Related work}
\lblsec{relatedwork}

{\bf Generative Adversarial Networks  (GANs)}~\cite{goodfellow2014generative,zhao2016energy} have achieved impressive results in image generation~\cite{denton2015deep,radford2015unsupervised}, image editing~\cite{zhu2016generative}, and representation learning~\cite{radford2015unsupervised,salimans2016improved,mathieu2016disentangling}.　
Recent methods adopt the same idea for conditional image generation applications, such as text2image~\cite{reed2016generative}, image inpainting~\cite{pathak2016context}, and future prediction~\cite{mathieu2015deep}, as well as to other domains like videos~\cite{vondrick2016generating} and 3D data~\cite{wu2016learning}. 
The key to GANs' success is the idea of an \textit{adversarial loss} that forces the generated images to be, in principle, indistinguishable from real photos.  
This loss is particularly powerful for image generation tasks,  
as this is exactly the objective that much of computer graphics aims to optimize.  We adopt an adversarial loss to learn the mapping such that the translated images cannot be distinguished from images in the target domain.

{\bf Image-to-Image Translation}
The idea of image-to-image translation goes back at least to Hertzmann et al.'s Image Analogies~\cite{hertzmann2001image}, who employ a non-parametric texture model~\cite{efros1999texture} on a single input-output training image pair.
More recent approaches use a \emph{dataset} of input-output examples to learn a parametric translation function using CNNs (e.g.,~\cite{long2015fully}). Our approach builds on the ``pix2pix" framework of Isola et al.~\cite{isola2016image}, which uses a conditional generative adversarial network~\cite{goodfellow2014generative} to learn a mapping from input to output images. Similar ideas have been applied to various tasks such as generating photographs from sketches~\cite{sangkloy2016scribbler} or from attribute and semantic layouts~\cite{karacan2016learning}. However, unlike the above prior work, we learn the mapping without paired training examples. 

{\bf Unpaired Image-to-Image Translation}
Several other methods also tackle the unpaired setting, where the goal is to relate two data domains: $X$ and $Y$. Rosales et al.~\shortcite{rosales2003unsupervised} propose a Bayesian framework that includes a prior based on a patch-based Markov random field computed from a source image and a likelihood term obtained from multiple style images. More recently, CoGAN~\cite{liu2016coupled} and cross-modal scene networks~\cite{aytar2016cross} use a weight-sharing strategy to learn a common representation across domains. 
Concurrent to our method, Liu et al.~\shortcite{liu2017unsupervised} extends the above framework with a combination of variational autoencoders~\cite{kingma2013auto} and generative adversarial networks~\cite{goodfellow2014generative}.
Another line of concurrent work~\cite{shrivastava2016learning,taigman2016unsupervised,bousmalis2016unsupervised} encourages the input and output to share specific ``content" features even though they may differ in ``style``. These methods also use adversarial networks, with additional terms to enforce the output to be close to the input in a predefined metric space, such as class label space~\cite{bousmalis2016unsupervised}, image pixel space~\cite{shrivastava2016learning}, and image feature space~\cite{taigman2016unsupervised}. 

Unlike the above approaches, our formulation does not rely on any task-specific, predefined similarity function between the input and output, nor do we assume that the input and output have to lie in the same low-dimensional embedding space. This makes our method a general-purpose solution for many vision and graphics tasks. We directly compare against several prior and contemporary approaches in \refsec{comparison}.

{\bf Cycle Consistency}
The idea of using transitivity as a way to regularize structured data has a long history.  In visual tracking, enforcing simple forward-backward consistency has been a standard trick for decades~\cite{kalal2010forward,sundaram2010dense}. In the language domain, verifying and improving translations via ``back translation and reconciliation'' is a technique used by human translators~\cite{brislin1970back} (including, humorously, by Mark Twain~\cite{twain1903}), as well as by machines~\cite{he2016dual}. 
More recently, higher-order cycle consistency has been used in
structure from motion~\cite{zach2010disambiguating},
3D shape matching~\cite{huang2013consistent}, co-segmentation~\cite{wang2013image}, dense semantic alignment~\cite{zhou2015flowweb,zhou2016learning}, and depth estimation~\cite{godard2016unsupervised}.  Of these, Zhou et al.~\cite{zhou2016learning} and Godard et al.~\cite{godard2016unsupervised} are most similar to our work, as they use a {\em cycle consistency loss} as a way of using transitivity to supervise CNN training.
In this work, we are introducing a similar loss to push $G$ and $F$ to be consistent with each other.  Concurrent with our work, in these same proceedings, Yi et al.~\cite{yi2017dualgan} independently use a similar objective for unpaired image-to-image translation, inspired by dual learning in machine translation~\cite{he2016dual}.

{\bf Neural Style Transfer}~\cite{gatys2015neural,johnson2016perceptual,ulyanov2016texture,gatys2016preserving} is another way to perform image-to-image translation, which synthesizes a novel image by combining the content of one image with the style of another image (typically a painting) based on matching the Gram matrix statistics of pre-trained deep features.  Our primary focus, on the other hand, is learning the mapping between two image collections, rather than between two specific images, by trying to capture correspondences between higher-level appearance structures. Therefore, our method can be applied to other tasks, such as painting$\rightarrow$ photo, object transfiguration, etc. where single sample transfer methods do not perform well. We compare these two methods in ~\refsec{application}.
